\documentclass[10pt,a4paper]{amsart}
\usepackage[margin=19mm]{geometry}
\usepackage{amsmath,amssymb,mathtools}
\usepackage[scr=rsfs,cal=euler]{mathalfa}
\usepackage{xfrac}
\usepackage[unicode,hidelinks,bookmarks=false]{hyperref}
\newcommand{\ie}{i.e.\ }
\newcommand{\cf}{cf.\ }
\newcommand{\resp}{respectively\ }
\newcommand{\Subsection}{\subsection}
\providecommand{\nobreakdash}{\nobreak\hbox{-}}
\setlength{\emergencystretch}{2em}
\multlinegap0pt
\usepackage{mathtools}
\usepackage{bm}
\usepackage{xcolor}
\usepackage{amssymb}
\usepackage{tikz}
\usepackage{algorithm}\usepackage{algpseudocode}
\newtheorem{theorem}{Theorem}
\newtheorem{lemma}{Lemma}
\usepackage{cleveref}
\crefname{theorem}{Theorem}{Theorems}
\crefname{lemma}{Lemma}{Lemmas}
\newcommand{\cD}{\mathcal{D}}
\newcommand{\cE}{\mathcal{E}}
\newcommand{\cH}{\mathcal{H}}
\newcommand{\cO}{\mathcal{O}}
\newcommand{\cS}{\mathcal{S}}
\newcommand{\codeg}{\Delta_2}
\newcommand{\ex}{\kappa}
\newcommand{\exx}{C_*}
\newcommand{\lam}{\lambda}
\newcommand{\tmix}{\tau_{\rm{mix}}}
\title{Sampling from the antiferromagnetic Ising model on bipartite, regular expander graphs}
\author{Anna Geisler, Mihyun Kang, Michail Sarantis, Ronen Wdowinski}
\date{Source: arXiv:2603.02101v1 (CC BY 4.0)}
\begin{document}
\maketitle

\noindent\emph{Source and licence.} Sampling from the antiferromagnetic Ising model on bipartite, regular expander graphs. Version arXiv:2603.02101v1 (CC BY 4.0), distributed by arXiv under the Creative Commons Attribution 4.0 International licence, http://creativecommons.org/licenses/by/4.0/, as recorded in the arXiv licence metadata for this exact version and shown on the abstract page https://arxiv.org/abs/2603.02101v1. Full licence notice: \texttt{licenses/arxiv-2603.02101-CC-BY-4.0.txt}.\par\smallskip

\noindent\textit{Editorial scope.} Counted unit: the article's theorem thm:slowmixing (source0.tex lines 275--280). The article's complete original proof of it is delivered with it (source0.tex lines 489--511, one \texttt{proof} environment of 23 lines, which the article places away from the statement). No mathematics is written, repaired or re-derived: the statement and the proof are the article's own lines, in the article's own notation, and no retained line is substituted or re-flowed.  Retained uncounted as the source-local context the proof needs: lines 380--385; lines 443--445; lines 465--470; lines 472--477; lines 479--484.  Nothing was omitted from any retained span, so the package carries no omission record.  No retained line contains a citation, so the wrapper carries no bibliography and no bibliography entry.  No retained line contains a figure environment, an image-inclusion command or drawing code, so no image is delivered and the package declares no asset dependency.  Editorial formatting only, in addition: an AMS \texttt{amsart} preamble that restates the article's own declarations and macros used by the retained lines, namely the environments \texttt{theorem}, \texttt{lemma} on separate counters, together with the standard packages those lines need.  The retained environments print in this document as Theorem 1, Lemma 1 in order of appearance, whereas the article prints the same items with the numbering of its own full document.  The complete native source of the article as distributed in the arXiv e-print for this exact version is bundled unchanged as \texttt{sources/195598/source0.tex}.
\begin{lemma} \label{lem:conduc_mix}
If $M$ is an ergodic Markov chain with conductance $\Phi$, then its mixing time satisfies
    \begin{equation*}
        \tmix(M)\geq\frac{1}{4\Phi}.
    \end{equation*}
\end{lemma}

\begin{equation} \label{eq:conductance_bound_weight_balanced}
    \Phi_M \leq \frac{\tilde{\omega}(\cS_{bal})}{(1+\lambda)^{n/2}}.
\end{equation}

\begin{lemma} \label{lem:small}
We have
\begin{align*}
    \tilde{\omega}(\cS_{small}) \leq (1+\lambda)^{n/4}.
\end{align*}
\end{lemma}

\begin{lemma} \label{lem:middle}
There exists a constant $C'>0$ such that
\begin{align*}
    \tilde{\omega}(\cS_{big, \cD}) \leq (1+\lambda)^{n/2} \exp\left( - \frac{C' \alpha^2 f(\lam)}{d^{\ex} \log d} \cdot n \right).
\end{align*}
\end{lemma}

\begin{lemma} \label{lem:rest}
We have
\begin{align*}
    \tilde{\omega}(\cS_{rest}) \leq (1+\lambda)^{n/2} \exp(-\Omega(\lam n)).
\end{align*}
\end{lemma}

\begin{theorem} \label{thm:slowmixing}
    Fix $\codeg \ge 1$, $0 \le \ex < 2$, and $\lam_0>0$. Then there exist constants $C_0, C > 0$ such that whenever $\lambda \le \lam_0$ and $\alpha \coloneqq \lam(1 - e^{-\beta}) \ge \frac{C_0 \log^{3/2} d}{d^{\exx}}$, the mixing time of the Glauber dynamics $M(G)$ on a graph $G \in \cH(n,d,\codeg,\ex)$ satisfies
    \[
    \tmix(M(G)) \geq \exp\left( \frac{C \alpha^2 f(\lam)}{d^{\ex} \log d} \cdot n\right).
    \]
\end{theorem}

\begin{proof}[Proof of \Cref{thm:slowmixing}]
Since $\cS_{bal} \subseteq \cS_{small} \cup \cS_{big, \cE} \cup \cS_{big, \cO} \cup \cS_{rest}$, we have 
    \[
    \tilde{\omega}(\cS_{bal}) \leq \tilde{\omega}(\cS_{small}) + \tilde{\omega}(\cS_{big, \cE}) + \tilde{\omega}(\cS_{big, \cO}) + \tilde{\omega}(\cS_{rest}).
    \]
    Plugging in \Cref{lem:small}, \Cref{lem:middle} and \Cref{lem:rest} yields that
    \begin{align} \label{eq:bound-sum}
        \tilde{\omega}(\cS_{bal}) \leq &(1+\lambda)^{n/2}\left[(1+\lam)^{-n/2}+2\exp\left( - \frac{C' \alpha^2 f(\lam)}{d^{\ex}\log d} \cdot n\right)+\exp(-\Omega(\lam n))\right]
        \end{align}
    We observe that the second summand is the dominating term.
    This follows from
    \[
    \frac{C' \alpha^2 f(\lam)}{d^{\ex}\log d} \cdot n \leq \frac{C' \lam_0^2}{R} \cdot \frac{1 }{d^{\ex}\log d} \cdot \lam n = o(\lam n) \hspace{10pt} \text{and} \hspace{10pt} \frac{C' \alpha^2 f(\lam)}{d^{\ex}\log d} \cdot n = o\left(\frac{n}{4} \log(1+\lam)\right),
    \]
    where the first upper bound uses $\alpha \leq \lam \leq \lam_0$, and the second bound uses $\log(1+\lambda) \geq \lam - \frac{\lam^2}{2}$. Thus, \eqref{eq:bound-sum} shows that there exists a constant $C > 0$ such that 
    \begin{align} \label{eq:weight_bal}
        \tilde{\omega}(\cS_{bal}) \leq (1+\lambda)^{n/2} \exp\left( - \frac{C \alpha^2 f(\lam) }{2 d^{\ex}\log d} \cdot n\right).
    \end{align}
    Plugging \eqref{eq:weight_bal} into \eqref{eq:conductance_bound_weight_balanced}, we conclude that the conductance of our Markov chain $M$ satisfies $\Phi_M \leq  \exp\left( - \frac{C \alpha^2 f(\lam)}{2 d^{\ex} \log d} \cdot n \right)$. Therefore, by \Cref{lem:conduc_mix}, the mixing time satisfies
    \[
    \tmix(M) \geq \exp\left( \frac{C \alpha^2 f(\lam)}{d^{\ex}\log d} \cdot n \right). \qedhere
    \]
\end{proof}

\end{document}
